\begin{helper}
Let \(V\) and \(V'\) be finite vertex sets, let \(\Delta\in\mathbb N\), and
let \(\gamma\in\mathbb R\).  Let
\[
  \nu \quad\text{on}\quad \mathcal P(V)\times\mathbb R^V,
  \qquad
  \nu' \quad\text{on}\quad \mathcal P(V')\times\mathbb R^{V'}
\]
be the source and target activation-priority product measures, with the
activation coordinate carrying the full measurable structure and the priority
coordinate carrying the finite product Borel structure.  Assume
\[
  0\le \gamma/\Delta\le 1.
\]
Assume that the source activation atoms satisfy, for every \(A\subseteq V\),
\[
  \nu\{(B,\pi):B=A\}
  =
  \operatorname{ofReal}\!\left(
    (\gamma/\Delta)^{|A|}
    (1-\gamma/\Delta)^{|V|-|A|}
  \right),
\]
and that source activation atoms have the lower-rectangle product law: for
every \(A\subseteq V\) and every threshold vector \(t:V\to\mathbb R\) with
\(t(v)\in[0,1]\) for all \(v\),
\[
  \nu\{(B,\pi):B=A,\ 0\le \pi(v)\le t(v)\text{ for every }v\in V\}
  =
  \nu\{(B,\pi):B=A\}\,
  \operatorname{ofReal}\!\left(\prod_{v\in V} t(v)\right).
\]
Assume the analogous two target laws: for every \(A'\subseteq V'\),
\[
  \nu'\{(B',\pi'):B'=A'\}
  =
  \operatorname{ofReal}\!\left(
    (\gamma/\Delta)^{|A'|}
    (1-\gamma/\Delta)^{|V'|-|A'|}
  \right),
\]
and, for every \(A'\subseteq V'\) and every \(t':V'\to\mathbb R\) with
\(t'(v')\in[0,1]\) for all \(v'\),
\[
  \nu'\{(B',\pi'):B'=A',\ 0\le \pi'(v')\le t'(v')
      \text{ for every }v'\in V'\}
  =
  \nu'\{(B',\pi'):B'=A'\}\,
  \operatorname{ofReal}\!\left(\prod_{v'\in V'} t'(v')\right).
\]
Let \(\phi:V\to V'\) be injective, let \(S\subseteq V\) be the source common
coordinate set, and let \(S'\subseteq V'\) be the target common coordinate
set with \(S'=\phi(S)\).  Then there is a paired measure
\(\lambda\) on
\[
  \bigl(\mathcal P(V)\times\mathbb R^V\bigr)
  \times
  \bigl(\mathcal P(V')\times\mathbb R^{V'}\bigr)
\]
such that the nonsynchronized common-record set has zero \(\lambda\)-mass,
\[
  \lambda\{((B,\pi),(B',\pi')):
     \neg(B'\cap S'=\phi(B\cap S)
       \text{ and }
       \pi'(\phi(v))=\pi(v)\text{ for every }v\in S)\}=0,
\]
and such that \(\lambda\)'s source and target projections are exactly
\(\nu\) and \(\nu'\) on all measurable events.  That is, for every
measurable source event \(E\subseteq\mathcal P(V)\times\mathbb R^V\) and
every measurable target event \(E'\subseteq\mathcal P(V')\times\mathbb R^{V'}\),
\[
  \lambda(\{(\eta,\eta'):\eta\in E\})=\nu(E),\qquad
  \lambda(\{(\eta,\eta'):\eta'\in E'\})=\nu'(E').
\]
The same $\lambda$ admits a
\noderef{SamplingSplitOutsideBlockerProductRealization} for
$p=\gamma/\Delta$, $\phi$, and $S$. Its entire joint law is the push-forward
of independent Bernoulli--uniform tagged coordinates by a measurable map.
\end{helper}
\begin{proof}
The tag construction below is supplied by \noderef{SamplingCopiedCommonTags};
its tag set $V\sqcup(V'\setminus\phi(S))$ is the same construction with
the common and source-only tags combined into the summand $V$.
The decoder is measurable by \noderef{SamplingCoordinateDecodeMeasurable}
on each projection. The two marginal cylinder calculations below are the
two applications of \noderef{SamplingInjectiveCoordinateDecodeLaws}, and
their upgrade to full marginal equalities is
\noderef{SamplingActivationRectangleMeasureUniqueness}. We spell out the
construction and the measure argument to identify the same joint witness.
Put $p=\gamma/\Delta$ and form the disjoint tag type
\[
T=S\sqcup((V\setminus S)\sqcup(V'\setminus S')).
\]
Send $v\in S$ to its common tag and $v\notin S$ to its source-only tag;
this defines $s$. Send $v'\in S'$ to the common tag of its unique preimage
under $\phi|_S$, and $v'\notin S'$ to its target-only tag; this defines $t$.
The preimage exists because $S'=\phi(S)$ and is unique by injectivity.
The disjoint sums make $s$ and $t$ injective. Every tag is used, and
$s(v)=t(v')$ holds exactly when $v\in S$ and $v'=\phi(v)$.

On each tag put $b_p\otimes u$, where $b_p=p\delta_1+(1-p)\delta_0$ on
$\mathbb R$ and $u=\operatorname{Leb}|_{[0,1]}$. Both are probability
measures, including $p=0$ and $p=1$. Let $m$ be their finite product over $T$.
Use $D=D_{s,t}$ from \noderef{SamplingSplitOutsideBlockerProductDecode}
and set $\lambda=D_*m$. The map $D$ is measurable: each priority coordinate
is a coordinate projection, and the inverse image of a prescribed source
activation set is a finite intersection of Borel sets
$\{q:q(s(v))_1=1\}$ and their complements. The target activation fibers
are treated identically. The activation spaces are finite and discrete, so
singleton fibers suffice; pairing these measurable maps gives the asserted
measurability into the natural product space. This supplies every field of
\noderef{SamplingSplitOutsideBlockerProductRealization}, including its
exact push-forward identity for this same $\lambda$.

For every input $q$, common activation and priority records of $D(q)$ agree
because they use the same common tag. Thus the nonsynchronized set has empty
inverse image. Each priority lies in $[0,1]$ almost surely, and the finite
union of exceptional coordinate sets is null. Consequently
\[
  \lambda(\text{nonsynchronized common records})=0,
\]
\[
  \lambda\{((B,\pi),(B',\pi')):\pi\notin[0,1]^V\}=0,\qquad
  \lambda\{((B,\pi),(B',\pi')):\pi'\notin[0,1]^{V'}\}=0,
\]
For a source atom $A$ and thresholds $a:V\to[0,1]$, the inverse cylinder
restricts the distinct tags $s(v)$. Each $v\in A$ contributes activation
factor $p$, each $v\notin A$ contributes $1-p$, and its priority factor
is $u([0,a(v)])=a(v)$. Tags outside $s(V)$ are unrestricted and contribute
one. Thus finite-product rectangle evaluation gives
$p^{|A|}(1-p)^{|V|-|A|}\prod_v a(v)$, also the value given by the
hypotheses on $\nu$. For a target atom $A'$, use the distinct tags $t(v')$;
the same calculation gives
$p^{|A'|}(1-p)^{|V'|-|A'|}\prod_{v'}a'(v')$, with source-only tags
integrating out. Renaming thresholds $t,t'$, we obtain
\[
  \lambda\{((B,\pi),(B',\pi')):B=A,\ 0\le\pi(v)\le t(v)\ \forall v\}
  =
  \nu\{(B,\pi):B=A,\ 0\le\pi(v)\le t(v)\ \forall v\},
\]
with the analogous target cylinder identity for \(V'\) and \(\nu'\).
The first of these conclusions is exactly the synchronization clause required
in the present statement.  It remains only to upgrade the cylinder identities
to full source and target marginal identities.

Write
\[
  Q_V=\{(B,\pi):0\le \pi(v)\le 1\hbox{ for every }v\in V\}.
\]
We first prove from the hypotheses on \(\nu\), and not from the paired
cylinder theorem, that \(\nu(Q_V^c)=0\).  Fix an activation atom
\[
  H_A=\{(B,\pi):B=A\}.
\]
Put \(t(v)=1\) for every \(v\).  The source lower-rectangle law gives
\[
  \nu(H_A\cap Q_V)
  =
  \nu(H_A)\operatorname{ofReal}\!\left(\prod_{v\in V}1\right)
  =
  \nu(H_A).
\]
The atom law identifies \(\nu(H_A)\) with a finite real Bernoulli mass, hence
\(\nu(H_A)<\infty\).  Since \(H_A\cap Q_V\subseteq H_A\), finite subtraction
in the measure algebra gives
\[
  \nu(H_A\setminus Q_V)=0.
\]
The finitely many atoms \(H_A\), \(A\subseteq V\), partition
\(\mathcal P(V)\times\mathbb R^V\).  Therefore
\[
  Q_V^c=\bigcup_{A\subseteq V}(H_A\setminus Q_V)
\]
up to equality, and finite subadditivity gives \(\nu(Q_V^c)=0\).

Let \(\mu_s\) be the source projection measure of \(\lambda\):
\[
  \mu_s(E)=\lambda\{((B,\pi),(B',\pi')):(B,\pi)\in E\}.
\]
The source unit-cube nullity proved above
says \(\mu_s(Q_V^c)=0\).  For every atom \(A\) the same all-one threshold
cylinder identity gives
\[
  \mu_s(H_A\cap Q_V)=\nu(H_A\cap Q_V)=\nu(H_A).
\]
Summing over the finite activation partition shows
\(\mu_s(Q_V)=\nu(Q_V)\).  Together with the two outside-cube nullities this
also gives \(\mu_s\) and \(\nu\) the same finite total mass.

We now prove equality on all measurable subsets of \(Q_V\).  On the trace
space \(Q_V\), consider the family
\[
  \mathcal P_s=
  \{H_A\cap\{(B,\pi):0\le \pi(v)\le t(v)\ \forall v\}:
    A\subseteq V,\ t(v)\in[0,1]\ \forall v\}.
\]
Adjoin the empty set to this family. It is then a \(\pi\)-system. Indeed,
the intersection of two activation
atoms is either empty or the same atom, and the intersection of two anchored
lower rectangles is the anchored lower rectangle with threshold
\(\min(t_1(v),t_2(v))\).  The empty set is harmless.  It generates the trace
product \(\sigma\)-algebra on \(Q_V\): the activation coordinate is a finite
discrete measurable space, and in each real coordinate the subspace Borel
\(\sigma\)-algebra on \([0,1]\) is generated by the anchored closed intervals
\([0,t]\), \(0\le t\le1\); finite products of these one-coordinate generators
generate the finite product Borel structure on \([0,1]^V\), and adjoining the
finite activation atoms gives the product with \(\mathcal P(V)\).

By the source cylinder identity computed above,
\(\mu_s\) and \(\nu\) agree on every member of \(\mathcal P_s\).  Since their
restrictions to \(Q_V\) have the same finite total mass, the usual
\(\pi\)-\(\lambda\) theorem applies: the class of trace-measurable
\(E\subseteq Q_V\) for which \(\mu_s(E)=\nu(E)\) is a Dynkin class containing
\(\mathcal P_s\), hence contains the trace \(\sigma\)-algebra generated by
\(\mathcal P_s\).  Thus \(\mu_s(E)=\nu(E)\) for every measurable
\(E\subseteq Q_V\).

Finally let \(E\subseteq \mathcal P(V)\times\mathbb R^V\) be any measurable
source event.  Decompose it as
\[
  E=(E\cap Q_V)\,\dot\cup\,(E\cap Q_V^c).
\]
The first piece is trace-measurable in \(Q_V\), so the preceding paragraph
gives equality of its \(\mu_s\)- and \(\nu\)-masses.  The second piece has
zero \(\mu_s\)-mass because \(\mu_s(Q_V^c)=0\), and zero \(\nu\)-mass because
\(\nu(Q_V^c)=0\).  Hence \(\mu_s(E)=\nu(E)\), which is the required source
projection identity for arbitrary measurable \(E\).

The target marginal is proved independently in exactly the same way, with
\[
  Q_{V'}=\{(B',\pi'):0\le \pi'(v')\le1\hbox{ for every }v'\in V'\}.
\]
The target atom law and target lower-rectangle law with all thresholds equal
to \(1\) first give \(\nu'(Q_{V'}^c)=0\) by finite additivity over the
activation atoms \(A'\subseteq V'\).  The target unit-cube nullity and target
cylinder identities computed above
give the corresponding projection measure
\[
  \mu_t(E')=\lambda\{((B,\pi),(B',\pi')):(B',\pi')\in E'\}
\]
zero mass outside \(Q_{V'}\) and equality with \(\nu'\) on all target
activation atoms times anchored lower rectangles.  Those target cylinders are
again a \(\pi\)-system generating the trace product \(\sigma\)-algebra on
\(Q_{V'}\).  The same finite-measure \(\pi\)-\(\lambda\) argument gives
\(\mu_t(E')=\nu'(E')\) for every measurable target event \(E'\), after
splitting \(E'\) into its part inside \(Q_{V'}\) and its null outside-cube
part.  Therefore the measure \(\lambda\) has the claimed full source and
target marginals, together with the already obtained almost-sure common
synchronization.
\end{proof}
